% Draft preprint. Every number comes from a file in results/; regenerate with `qudock reproduce`.
% Build: cd paper && latexmk -pdf qudock.tex
\documentclass[11pt]{article}
\usepackage[margin=1in]{geometry}
\usepackage{amsmath,amssymb}
\usepackage{booktabs}
\usepackage{graphicx}
\usepackage[hidelinks]{hyperref}
\usepackage[numbers]{natbib}

\title{Do quantum kernels help find Alzheimer's drug candidates?\\
\large A scaffold-split benchmark of quantum and classical models for BACE-1 inhibition,\\
with measurements on IBM Quantum hardware}
\author{Shubhodeep Chatterjee \and Srijita Banerjee \and Rounak Pal\\
\small Team BlackOut, VTU-Qiskit Fall Fest 2026}
\date{}

\begin{document}
\maketitle

\begin{abstract}
Quantum kernel methods are often proposed for drug discovery on the grounds that they may learn
more from fewer labelled molecules. We test that claim on inhibition of $\beta$-secretase~1
(BACE-1), an Alzheimer's disease target, under conditions designed to make a fair comparison:
Bemis--Murcko scaffold splits, identical features for quantum and classical models,
hyperparameters tuned by inner cross-validation, multiple seeds and paired statistics. We
evaluate the ZZ fidelity kernel on unsupervised (PCA) and supervised (PLS) angles with a
bandwidth tuned by kernel-target alignment, the projected quantum kernel, a feature map whose
own parameters are trained, a hybrid quantum + Tanimoto kernel, and a new feature map whose
entanglement follows each molecule's pharmacophore graph. We measure kernels on simulated noisy
devices and on IBM's \texttt{ibm\_marrakesh}. Better features and tuned bandwidths lift the
quantum kernel substantially (ROC-AUC 0.708 to 0.743 at full size), and a hybrid kernel comes
within 0.006 of the best classical model, but no quantum model beats the strongest classical
baseline on BACE-1 or on blood-brain barrier penetration (BBBP). Hardware noise distorts the
kernel by 17.5\% yet preserves its shape (correlation 0.996 with the exact kernel). We release
the pipeline, a multi-objective screening tool and QuDock-Bench, a fixed protocol for scoring
any molecule classifier on BACE, BBBP and ClinTox.
\end{abstract}

\section{Introduction}
Quantum kernels embed data in the Hilbert space of a parametrised circuit and compare points by
state overlap~\cite{havlicek2019}. Their appeal for chemistry rests on the hope that such
embeddings capture structure classical kernels miss, and that this pays off when labelled data
are scarce. Testing that hope requires care: random splits leak near-duplicate molecules
between training and test sets, quantum models are often compared against untuned classical
baselines, and results on a single seed are noisy~\cite{bowles2024}.

We ask a narrow question carefully: on BACE-1 inhibition~\cite{wu2018}, does any quantum kernel
we can build, simulate and run outperform well-tuned classical models on unseen chemical
scaffolds? Our contributions are
(i) a fair comparison across five families of quantum kernels and six classical baselines,
(ii) a pharmacophore-graph feature map in which entanglement encodes molecular structure,
(iii) noise studies on simulated and real IBM devices,
(iv) calibrated confidence and atom-level explanations for the quantum models, and
(v) QuDock-Bench, a reusable benchmark with a plugin interface.

\section{Methods}
\paragraph{Data.} MoleculeNet BACE (1{,}513 molecules; 1{,}210 train and 303 test after a
scaffold split, rarest scaffolds in the test set). SMILES are canonicalised and molecules with
conflicting labels removed. BBBP (1{,}965 molecules) and ClinTox serve as second and third tasks.

\paragraph{Features.} ECFP4 fingerprints (2{,}048 bits), optionally with ten RDKit descriptors,
are standardised and compressed to one value per qubit by PCA or by PLS fit on training labels,
then rescaled to $[0,\pi]$.

\paragraph{Quantum kernels.}
\emph{ZZ fidelity kernel}: $k(x,y)=|\langle\phi(x)|\phi(y)\rangle|^2$ for the ZZ feature map with
linear entanglement and bandwidth $b$ scaling the angles~\cite{shaydulin2022}; $b$ is chosen by
kernel-target alignment on training data.
\emph{Projected kernel}~\cite{huang2021}: an RBF kernel on single-qubit Pauli expectations.
\emph{Trainable map}~\cite{hubregtsen2022}: per-qubit encoding weights, per-pair entangling
weights and RY rotations, trained by alignment from the ZZ map as starting point.
\emph{Hybrid}: $w\,k_Q + (1-w)\,k_{\text{Tanimoto}}$ with $w$ chosen by inner CV.
\emph{Pharmacophore-graph map} (this work): six qubits, one per pharmacophore family (donor,
acceptor, positive and negative ionizable, aromatic, hydrophobe). Qubit $q$ is rotated by
$x_q=\pi c_q/(c_q+2)$ for $c_q$ features of that family, and every pair is entangled by a ZZ
rotation of strength $w_{qr}=\pi s_{qr}/(s_{qr}+1)$, where
$s_{qr}=\sum \exp(-d/2)$ sums over feature pairs at topological distance $d$. The entanglement
pattern therefore differs per molecule and follows its pharmacophore geometry.

\paragraph{Scaling tools.} A Nystr\"om map $\phi(x)=k(x,L)K_{LL}^{-1/2}$ with $m$ landmarks reduces
kernel evaluations from $n^2/2$ to $nm$; classical shadows~\cite{huang2020} estimate all one-
and two-qubit Pauli expectations from the same randomised measurements.

\paragraph{Evaluation.} Test ROC-AUC, PR-AUC, enrichment at 1\% and 5\%, and BEDROC. Every model
sees the same molecules; training subsets are stratified and seeded; differences are tested by
paired bootstrap on the test set or Wilcoxon tests over seeds.

\section{Results}

\begin{table}[h]
\centering
\caption{Mean test ROC-AUC over 10 seeds on BACE (6 qubits).}
\label{tab:final}
\small
\begin{tabular}{rcccccc}
\toprule
$n_\text{train}$ & QSVC PCA & QSVC PLS & SVM-RBF PLS & Hybrid PCA & Tanimoto SVM & RF fingerprint\\
\midrule
50 & 0.607 & 0.656 & 0.675 & 0.651 & 0.655 & 0.684\\
100 & 0.647 & 0.687 & 0.710 & 0.702 & 0.712 & 0.735\\
200 & 0.663 & 0.712 & 0.737 & 0.726 & 0.742 & 0.756\\
400 & 0.693 & 0.732 & 0.742 & 0.757 & 0.762 & 0.769\\
1{,}210 & 0.708 & 0.743 & 0.769 & 0.776 & 0.781 & 0.782\\
\bottomrule
\end{tabular}
\end{table}

\paragraph{Better features help the quantum kernel; the quantum kernel does not beat its
classical twin.} Supervised PLS angles and an alignment-tuned bandwidth raise QSVC by 0.035 to
0.049 ROC-AUC at every training size (Table~\ref{tab:final}). On the very same PLS features,
however, SVM-RBF remains ahead by 0.026 at full size (paired bootstrap $p=0.03$). The hybrid
kernel keeps a non-zero quantum weight under inner CV and reaches 0.776, statistically level
with the Tanimoto SVM ($p=0.27$) but not above it. On BBBP the picture repeats: hybrid 0.827
against 0.826 (Tanimoto SVM) and 0.829 (Random Forest).

\paragraph{No small-data advantage.} With 25 or 50 training molecules the quantum kernel is the
weakest model in the first-round benchmark, contrary to the motivating hypothesis.

\paragraph{Concentration is real and manageable.} With a fixed bandwidth the spread of
off-diagonal kernel values falls 11-fold from 4 to 14 qubits; alignment-tuned bandwidths keep it
9 to 17 times larger and recover 0.03 to 0.09 ROC-AUC at every width.

\paragraph{Training the feature map helps a little.} Alignment-trained maps win 14 of 20 runs,
+0.012 ROC-AUC on average (Wilcoxon $p=0.03$), with the largest gain at 100 molecules (+0.021).

\paragraph{Encoding chemistry in the entanglement does not add signal.} The pharmacophore-graph
kernel ties an RBF SVM on exactly the same 21 numbers (0.673 vs 0.666 at 400 molecules; 0.656 vs
0.682 at full size) and loses to a Random Forest on them (0.745).

\paragraph{Hardware.} On \texttt{ibm\_marrakesh} (4 qubits, 12 molecules, 66 circuits, 1{,}024
shots) the measured kernel deviates from the exact kernel by 17.5\% (FakeTorino simulation:
10.3\%) yet correlates at 0.996, so noise mostly shrinks kernel values uniformly and leaves the
ranking intact. In simulation, device noise contributes 8--13\% kernel error that more shots do
not remove, while test ROC-AUC stays within 0.62--0.66 across all settings.

\paragraph{Confidence and explanations.} Conformal prediction calibrated on rare training
scaffolds delivers 66--86\% coverage when 90\% is requested under a scaffold split, an honest
measure of how new the test chemistry is. Occlusion-based atom attributions are faithful:
removing the highlighted fifth of atoms costs 0.563 score against 0.381 for random atoms while
breaking fewer fingerprint bits.

\paragraph{Multi-objective screening.} Ranking the BACE test set by activity alone gives ROC-AUC
0.739; ranking by a geometric mean of calibrated activity, predicted brain penetration and QED
gives 0.703, the expected cost of selecting for molecules that can also reach the brain.

\paragraph{QuDock-Bench.} Table~\ref{tab:bench} scores seven models on three tasks under one
protocol. Classical models take the top three places on every task; the best quantum entry,
the hybrid kernel, ranks fourth. The gap is largest on ClinTox, where only 6\% of molecules are
positive and the quantum kernels on six compressed angles fall to 0.52--0.69 while fingerprint
models reach 0.84--0.87. At full training size the seeds only change cross-validation folds and
model randomness, so several intervals are close to zero. The Random Forest and gradient
boosting entries use balanced class weights and the SVM entries do not; the unbalanced Tanimoto
SVM still ranks first.

% QuDock-Bench leaderboard: regenerate with python paper/make_tables.py
\begin{table}[h]
\centering
\caption{QuDock-Bench: test ROC-AUC (mean $\pm$ 95\% interval over 3 seeds) on scaffold splits, full training sets.}
\label{tab:bench}
\small
\begin{tabular}{lcccc}
\toprule
Model & BACE & BBBP & CLINTOX & Mean rank\\
\midrule
Tanimoto SVM & 0.784 $\pm$ 0.008 & 0.826 $\pm$ 0.000 & 0.865 $\pm$ 0.000 & 1.7\\
Gradient boosting, ECFP4 + descriptors & 0.759 $\pm$ 0.007 & 0.831 $\pm$ 0.006 & 0.850 $\pm$ 0.010 & 2.3\\
Random Forest, ECFP4 & 0.782 $\pm$ 0.002 & 0.831 $\pm$ 0.001 & 0.842 $\pm$ 0.007 & 2.3\\
Hybrid ZZ + Tanimoto (quantum) & 0.766 $\pm$ 0.000 & 0.784 $\pm$ 0.000 & 0.691 $\pm$ 0.010 & 4.3\\
QSVC, PLS (quantum) & 0.754 $\pm$ 0.000 & 0.778 $\pm$ 0.000 & 0.622 $\pm$ 0.021 & 5.7\\
Projected kernel, Nystr\"om (quantum) & 0.716 $\pm$ 0.004 & 0.795 $\pm$ 0.006 & 0.518 $\pm$ 0.007 & 5.7\\
Pharmacophore-graph map (quantum) & 0.656 $\pm$ 0.000 & 0.789 $\pm$ 0.008 & 0.579 $\pm$ 0.073 & 6.0\\
\bottomrule
\end{tabular}
\end{table}


\section{Discussion}
Across five families of quantum kernels, two datasets, simulated and real hardware, we find the
same pattern: quantum kernels become competitive when given good features and tuned
bandwidths, and hybrids with a classical kernel come within statistical noise of the best
classical model, but none surpasses it. The gains we measured come from features, bandwidth
selection and kernel training, all of which have classical analogues. We consider this a useful
negative result, and we release the benchmark so that future quantum feature maps can be judged
on the same footing.

\paragraph{Limitations.} Six to fourteen qubits; one hardware run of 12 molecules for the kernel
study; no graph neural network baseline; PLS compression sees training labels, which makes
inner-CV estimates of PLS-based models slightly optimistic.

\paragraph{Code and data.} \url{https://github.com/vtu-qiskit-fall-fest-2026/vtu-hc-blackout}.
Every table is regenerated by \texttt{qudock reproduce}; each results folder carries a
provenance record with the code version, package versions and data hashes.

\bibliographystyle{unsrtnat}
\bibliography{refs}
\end{document}
